% =============================================================================
%  §1  Ch.4 最优耦合的基本性质
%  来源：Villani, Optimal Transport: Old and New, Chapter 4 (pp. 43–48)
% =============================================================================
\TLsection{1 \ 最优耦合的基本性质}

% ── 动机：为什么存在性不是显然的 ───────────────────────────────────────────────
\begin{frame}{为什么最优耦合"存在"不是显然的？}
  我们已经知道 Kantorovich 问题的可行集 $\coup{\mu}{\nu}$ 非空（含乘积测度 $\mu\otimes\nu$），
  目标 $I[\pi]=\int_{\X\times\Y} c\dd\pi$ 对 $\pi$ 线性。

  \textbf{朴素想法。} 取一列"趋向下确界"的计划：存在 $(\pi_k)\subset\coup{\mu}{\nu}$ 使
  \[
    I[\pi_k] \;\To\; \inf_{\pi\in\coup{\mu}{\nu}} I[\pi].
  \]
  \textbf{问题。} $(\pi_k)$ 是否有收敛子列？极限是否还在 $\coup{\mu}{\nu}$ 里？

  \begin{columns}[T]
    \begin{column}{0.5\textwidth}
      \textbf{两个关键工具：}
      \begin{itemize}
        \item \textbf{紧性。} $\coup{\mu}{\nu}$ 在弱拓扑下紧——使得收敛子列存在。
        \item \textbf{下半连续。} $I[\pi]$ 在弱极限处不增——使得极限是最优。
      \end{itemize}
    \end{column}
    \begin{column}{0.5\textwidth}
      \textbf{对应结果：}
      \begin{itemize}
        \item Lemma 4.4 $+$ Prokhorov 定理 $\Rightarrow$ 紧性。
        \item Lemma 4.3 $\Rightarrow$ 下半连续。
        \item Theorem 4.1 组装：最优耦合存在。
      \end{itemize}
    \end{column}
  \end{columns}
  \TLtakeaway{存在性证明的骨架：下确界列 $\to$ 紧性（取极限） $\to$ 下半连续（极限仍最优）。}
\end{frame}

% ── 定理 4.1 陈述 ─────────────────────────────────────────────────────────────
\begin{frame}{定理 4.1：最优耦合的存在性}
  \deflab{定理 4.1（最优耦合存在，Villani Thm 4.1）。}

  设 $\X,\Y$ 为 Polish 空间，$\mu\in\Pp(\X)$，$\nu\in\Pp(\Y)$。
  设 $a:\X\to\R\cup\{-\infty\}$ 上半连续且 $a\in L^1(\mu)$，
  $b:\Y\to\R\cup\{-\infty\}$ 上半连续且 $b\in L^1(\nu)$。
  设成本函数 $c:\X\times\Y\to\R\cup\{+\infty\}$ 下半连续，且
  \[
    c(x,y) \;\ge\; a(x)+b(y) \quad \forall\, x\in\X,\; y\in\Y.
  \]
  则\emph{存在}一个耦合 $\pi^\star\in\coup{\mu}{\nu}$ 使得
  \[
    I[\pi^\star] = \inf_{\pi\in\coup{\mu}{\nu}} I[\pi], \quad I[\pi]:=\int_{\X\times\Y} c(x,y)\dd\pi(x,y).
  \]

  \deflab{注记 4.2。} 下界条件 $c\ge a+b$ 保证 $I[\pi]$ 在 $\R\cup\{+\infty\}$ 中\emph{有定义}（不出现 $+\infty-\infty$）。
  多数应用中可取 $a=b=0$，即 $c\ge0$。

  \TLtakeaway{定理 4.1 是"第一件好事"：在 Polish 空间上、下半连续成本满足下界条件时，最优耦合必然存在。}
\end{frame}

% ── 注记 4.5：存在性 ≠ 有限性 ────────────────────────────────────────────────
\begin{frame}{注记 4.5：存在性不等于有限最优成本}
  \deflab{注记 4.5。} 定理 4.1 断言最优耦合\emph{存在}，但不排除
  \[
    \inf_{\pi\in\coup{\mu}{\nu}} I[\pi] = +\infty.
  \]
  即：存在最优计划，但其成本可能无穷大。

  \textbf{何时成本有限？} 一个简单的充分条件：
  \[
    \iint_{\X\times\Y} c(x,y)\dd\mu(x)\dd\nu(y) < +\infty.
  \]
  此条件等价于\emph{独立耦合} $\mu\otimes\nu$ 的成本有限，从而保证至少有一个计划成本有限，
  进而最优成本有限。

  \textbf{更强的条件（后续章节用到）：} 若
  $c(x,y)\le c_{\X}(x)+c_{\Y}(y)$，其中 $c_{\X}\in L^1(\mu)$，$c_{\Y}\in L^1(\nu)$，
  则\emph{任意}耦合的成本均有限。

  \TLtakeaway{成本有限性是独立的条件：验证 $\mu\otimes\nu$ 的成本是否有限，是最简单的途径。}
\end{frame}

% ── Prokhorov 定理回顾 + 引理 4.4 ────────────────────────────────────────────
\begin{frame}{Prokhorov 定理与耦合集的胶着性（引理 4.4）}
  \deflab{Prokhorov 定理（紧性判据）。}
  设 $\X$ 为 Polish 空间，$\mathcal{F}\subset\Pp(\X)$。
  $\mathcal{F}$ 在弱拓扑下\emph{预紧}当且仅当它\TLterm{胶着 (tight)}：
  \[
    \forall\eps>0,\;\exists\text{ 紧集 } K_\eps\subset\X, \;
    \mu[\X\setminus K_\eps]\le\eps \;\forall\,\mu\in\mathcal{F}.
  \]
  直觉：无质量"逃逸到无穷"。

  \deflab{引理 4.4（转移计划的胶着性）。}
  设 $\mathcal{F}\subset\Pp(\X)$ 与 $\mathcal{G}\subset\Pp(\Y)$ 均胶着，
  则 $\Pi(\mathcal{F},\mathcal{G})$ 在 $\Pp(\X\times\Y)$ 中\emph{也胶着}。

  \textbf{证明思路。} 对任意 $\eps>0$，取紧集 $K_\eps\subset\X$，$L_\eps\subset\Y$ 使得
  $\mu[\X\setminus K_\eps]\le\eps$，$\nu[\Y\setminus L_\eps]\le\eps$。
  对任意耦合 $(X,Y)\sim\pi\in\Pi(\mathcal{F},\mathcal{G})$：
  \[
    \Prob\bigl[(X,Y)\notin K_\eps\times L_\eps\bigr]
    \;\le\; \Prob[X\notin K_\eps]+\Prob[Y\notin L_\eps]
    \;\le\; 2\eps.
  \]
  $K_\eps\times L_\eps$ 在 $\X\times\Y$ 中紧，上界 $2\eps$ 与耦合选取无关。\hfill$\square$

  \TLtakeaway{单个 $\mu,\nu$ 的胶着性（在 Polish 空间上自动成立）可"提升"到整个耦合集 $\coup{\mu}{\nu}$。}
\end{frame}

% ── 引理 4.3：成本泛函的下半连续性 ──────────────────────────────────────────
\begin{frame}{引理 4.3：成本泛函的下半连续性（陈述）}
  \deflab{引理 4.3（成本泛函的下半连续性）。}
  设 $\X,\Y$ 为 Polish 空间，$c:\X\times\Y\to\R\cup\{+\infty\}$ 下半连续，
  $h:\X\times\Y\to\R\cup\{-\infty\}$ 上半连续且 $c\ge h$。
  设 $\pi_k\wkto\pi$，且 $h\in L^1(\pi_k)$，$h\in L^1(\pi)$，
  $\int h\dd\pi_k\To\int h\dd\pi$。则
  \[
    \int_{\X\times\Y} c\dd\pi \;\le\; \liminf_{k\to\infty}\int_{\X\times\Y} c\dd\pi_k.
  \]

  \textbf{读法。} 泛函 $\pi\mapsto\int c\dd\pi$ 对弱收敛\emph{下半连续}——这正是 §0 中"下半连续"概念在测度空间 $\Pp(\X\times\Y)$ 上的体现（辅助函数 $h$ 仅用于把成本下界 $c\ge a+b$ 纳入框架）。

  \TLtakeaway{下半连续成本泛函对弱收敛也下半连续——这让"取极限不增加成本"成为可能（定理 4.1 的另一半引擎）。}
\end{frame}

% ── 引理 4.3 的证明 ──────────────────────────────────────────────────────────
\begin{frame}{引理 4.3 的证明（一）：为什么可令 $c\ge0$}
  \textbf{先解释归约。}\proofskipnote 因 $c\ge h$，函数 $c-h$ 非负；若已知引理对非负成本 $c-h$ 成立，则
  \[
    \begin{aligned}
      \int c\dd\pi
      &= \int(c-h)\dd\pi+\int h\dd\pi\\
      &\le \liminf_{k}\int(c-h)\dd\pi_k+\lim_{k}\int h\dd\pi_k\\
      &= \liminf_{k}\int c\dd\pi_k .
    \end{aligned}
  \]
  最后一行用的是假设 $\int h\dd\pi_k\To\int h\dd\pi$。所以只需证明非负下半连续情形。

  \textbf{合法性。} $h$ 的积分沿 $\pi_k\wkto\pi$ 收敛，表示辅助下界项不会在极限中丢失；
  因此给 $c-h$ 证明的下半连续不等式，可以通过加回同一个收敛的 $h$ 项传回原成本 $c$。

  \TLtakeaway{归约不是省略：$c\ge h$ 给出非负成本 $c-h$，而 $\int h\dd\pi_k\to\int h\dd\pi$ 保证加回 $h$ 后不等式仍成立。}
\end{frame}

\begin{frame}{引理 4.3 的证明：逼近 + 单调收敛}
  \deflab{补一个事实：单调收敛定理。}
  若 $0\le c_{\ell}\nearrow c$，则 $\sup_{\ell}\int c_{\ell}\dd\pi=\lim_{\ell}\int c_{\ell}\dd\pi=\int c\dd\pi$。\proofskipnote

  \textbf{非负情形。} 用 §0 下半连续性质 (a)（inf-卷积）取有界连续逼近
  $c_{\ell}\in C_b(\X\times\Y)$，$0\le c_{\ell}\nearrow c$。固定 $\ell$，弱收敛给
  $\int c_{\ell}\dd\pi_k\to\int c_{\ell}\dd\pi$，且 $c_{\ell}\le c$，故
  \[
    \int c_{\ell}\dd\pi
    = \lim_{k}\int c_{\ell}\dd\pi_k
    \le \liminf_{k}\int c\dd\pi_k
    \quad\Longrightarrow\quad
    \int c\dd\pi\le \liminf_{k}\int c\dd\pi_k .
  \]
  右端与 $\ell$ 无关，对 $\ell$ 取上确界，并用上面的单调收敛事实，即得结论。\hfill$\square$
  \TLtakeaway{非负情形的核心：lsc 给出连续递增逼近；固定 $\ell$ 用弱收敛和 $c_{\ell}\le c$，最后用单调收敛回到 $c$。}
\end{frame}

% ── 定理 4.1 的证明（一a）：预紧 ─────────────────────────────────────────────
\begin{frame}{定理 4.1 的证明（一a）：$\coup{\mu}{\nu}$ 为何\emph{预紧}}
  \textbf{总思路。}\proofskipnote 把 §0 的"紧集上 lsc 必达最小"搬到测度空间：取"紧集"$=\coup{\mu}{\nu}$、"lsc 目标"$=I[\pi]=\int_{\X\times\Y} c\dd\pi$。本帧先证可行集\emph{预紧}（弱闭包紧）。

  \deflab{补一个事实：Ulam 定理。} Polish 空间上\emph{每个}概率测度 $\mu$ 都自动胶着：完备可分度量空间里 $\mu$ 对紧集\TLterm{内正则 (inner regular)}——任给 $\eps>0$ 总有紧集 $K_\eps$ 使 $\mu[\X\setminus K_\eps]\le\eps$。故单点族 $\{\mu\},\{\nu\}$ 胶着。

  \textbf{推出预紧（三步，每步标注依据）：}
  \begin{itemize}\small
    \item $\{\mu\},\{\nu\}$ 胶着\hfill（Ulam 定理）
    \item $\Rightarrow$ $\coup{\mu}{\nu}$ 在 $\Pp(\X\times\Y)$ 中胶着\hfill（引理 4.4：两端胶着"提升"到耦合集）
    \item $\Rightarrow$ $\coup{\mu}{\nu}$ \emph{预紧}\hfill（Prokhorov：胶着 $\Leftrightarrow$ 弱预紧）
  \end{itemize}

  \TLtakeaway{预紧来自一条链：Ulam（单个测度胶着）$\to$ 引理 4.4（耦合集胶着）$\to$ Prokhorov（胶着即预紧）。}
\end{frame}

% ── 定理 4.1 的证明（一b）：闭性 → 紧，取极小化列 ────────────────────────────
\begin{frame}{定理 4.1 的证明（一b）：从预紧到\emph{紧}，并取极小化列}
  \textbf{预紧还不够：}极限未必落在 $\coup{\mu}{\nu}$ 里，故须再证它\emph{闭}（闭 $+$ 预紧 $=$ 紧）。\proofskipnote

  \deflab{闭性：把"边际是 $\mu$"传给极限。}
  设 $\pi_k\in\coup{\mu}{\nu}$，$\pi_k\wkto\pi$。任取 $\varphi\in C_b(\X)$，视 $(x,y)\mapsto\varphi(x)$ 为有界连续函数：
  \begin{align*}
    \int_{\X}\varphi\dd\big(\pf{(\mathrm{proj}_{\X})}{\pi}\big)
    &=\int_{\X\times\Y}\!\varphi(x)\dd\pi
    \;\overset{\pi_k\wkto\pi}{=}\;\lim_{k}\int_{\X\times\Y}\!\varphi(x)\dd\pi_k\\
    &=\lim_{k}\int_{\X}\varphi\dd\mu=\int_{\X}\varphi\dd\mu .
  \end{align*}
  对一切 $\varphi$ 成立 $\Rightarrow$ $\pf{(\mathrm{proj}_{\X})}{\pi}=\mu$，同理第二边际 $=\nu$，故 $\pi\in\coup{\mu}{\nu}$：集合\emph{闭}。\textbf{闭 $+$ 预紧 $\Rightarrow$ $\coup{\mu}{\nu}$ 弱紧。}

  \TLtakeaway{闭性核心只有一行：用检验函数 $\varphi(x)$ 测两边，弱收敛把"边际 $=\mu$"原样传到极限；再配预紧即得\emph{紧}。}
\end{frame}

% ── 定理 4.1 的证明（二）：下半连续锁定最优性 ────────────────────────────────
\begin{frame}{定理 4.1 的证明（二）：用下半连续\emph{锁定}最优性}
  \textbf{取极小化列。}\proofskipnote 取 $(\pi_k)$ 使 $I[\pi_k]\to\inf I$，（一b）弱紧性抽出子列 $\pi_k\wkto\pi^\star\in\coup{\mu}{\nu}$。难点：弱收敛\emph{不}保证 $I[\pi_k]\to I[\pi^\star]$；但 lsc（引理 4.3）给的\emph{一个方向}够用。

  \deflab{先核对引理 4.3 的前提。} 取 $h(x,y)=a(x)+b(y)$（下界条件 $c\ge a+b$ 中的 $h$，上半连续）。\textbf{关键：}所有 $\pi_k,\pi^\star$ 共享边际 $\mu,\nu$，故
  \[
    \int h\dd\pi_k=\int_{\X} a\dd\mu+\int_{\Y} b\dd\nu=\int h\dd\pi^\star .
  \]
  右端与 $k$ 无关（\emph{常数列}），故前提"$\int h\dd\pi_k\to\int h\dd\pi^\star$"自动成立，引理 4.3 适用。

  \textbf{收尾（夹逼）：}
  \begin{itemize}\small
    \item 引理 4.3 $\Rightarrow$ $I[\pi^\star]\le\liminf_k I[\pi_k]=\inf I$；
    \item 可行性 $\Rightarrow$ $I[\pi^\star]\ge\inf I$；夹逼得 $I[\pi^\star]=\inf I$，即 $\pi^\star$ \emph{最优}。\hfill$\square$
  \end{itemize}

  \TLtakeaway{lsc 只需"成本不会向上跳"这一个方向 $I[\pi^\star]\le\liminf I[\pi_k]=\inf I$，配可行性即夹出等号。}
\end{frame}

% ── 定理 4.6：最优性被子计划继承 ─────────────────────────────────────────────
\begin{frame}{第二件好事：最优性被子计划继承（定理 4.6）}
  \deflab{定理 4.6（限制性质，Villani Thm 4.6）。}
  设定同定理 4.1，并设最优传输成本 $C(\mu,\nu)<+\infty$，$\pi\in\coup{\mu}{\nu}$ 为最优计划。
  设 $\pi'$ 是 $\X\times\Y$ 上的非负测度，满足 $\pi'\le\pi$（作为测度）且 $\pi'[\X\times\Y]>0$。
  定义归一化限制：
  \[
    \tilde{\pi} \;:=\; \frac{\pi'}{\pi'[\X\times\Y]}.
  \]
  则 $\tilde{\pi}$ 是其边际 $\tilde{\mu}=\pf{(\mathrm{proj}_{\X})}{\tilde{\pi}}$，
  $\tilde{\nu}=\pf{(\mathrm{proj}_{\Y})}{\tilde{\pi}}$ 之间的\emph{最优}耦合。

  \textbf{通俗理解。} "最优传输计划的任意一个子计划，对其所搬运的那部分质量来说，也是最优的。"
  否则可以局部改进子计划，从而降低整体成本——矛盾。

  \TLtakeaway{限制性质说明：最优计划\emph{局部也优}。这是全书反复使用的简单而有力的原则。}
\end{frame}

% ── 定理 4.6 的矛盾证明（一）：构造与验证 ────────────────────────────────────
\begin{frame}{定理 4.6 的证明（一）：构造竞争计划}
  \textbf{反证。}\proofskipnote 令 $Z=\pi'[\X\times\Y]>0$，$\tilde{\pi}=\pi'/Z$。若 $\tilde{\pi}$ 不是最优的，
  则存在同边际计划 $\hat{\pi}\in\coup{\tilde{\mu}}{\tilde{\nu}}$，满足 $I[\hat{\pi}]<I[\tilde{\pi}]$。

  \deflab{补一个事实：同边际计划的差有零边际。}
  因 $\hat{\pi}$ 与 $\tilde{\pi}$ 共享边际 $(\tilde{\mu},\tilde{\nu})$，所以
  \[
    \pf{(\mathrm{proj}_{\X})}{(\hat{\pi}-\tilde{\pi})}=0,\quad
    \pf{(\mathrm{proj}_{\Y})}{(\hat{\pi}-\tilde{\pi})}=0,\qquad
    \hat{\pi}_{\mathrm{new}}:=(\pi-\pi')+Z\hat{\pi}=\pi+Z(\hat{\pi}-\tilde{\pi}).
  \]
  这里 $\pi-\pi'$ 非负（因为 $\pi'\le\pi$），$Z\hat{\pi}$ 非负，故 $\hat{\pi}_{\mathrm{new}}$ 是非负测度。
  又因上式右端的差项有零边际，$\hat{\pi}_{\mathrm{new}}$ 的两个边际仍是 $\mu,\nu$，即它是 $\pi$ 的竞争计划。

  \textbf{成本严格下降。} $C(\mu,\nu)<+\infty$ 保证下面的加减运算合法；于是
  \[
    I[\hat{\pi}_{\mathrm{new}}]
    = I[\pi]-Z I[\tilde{\pi}]+Z I[\hat{\pi}]
    < I[\pi].
  \]
  这与 $\pi$ 最优性矛盾，故 $\tilde{\pi}$ 是最优的。\hfill$\square$
  \TLtakeaway{隐藏步骤是边际核对：把新计划写成 $\pi+Z(\hat{\pi}-\tilde{\pi})$，同边际差项的边际为零，所以替换只改变局部成本，不改变全局边际。}
\end{frame}

% ── 定理 4.6 的证明（二）：唯一性的传递 ─────────────────────────────────────
\begin{frame}{定理 4.6 的证明（二）：唯一性的传递}
  \textbf{目标。}\proofskipnote 已知 $\pi$ 是 $(\mu,\nu)$ 的\emph{唯一}最优计划；要证 $\tilde{\pi}$ 是
  $(\tilde{\mu},\tilde{\nu})$ 的唯一最优计划。
  \textbf{逐步传递唯一性。}
  \begin{itemize}\small
    \item 任取一个 $(\tilde{\mu},\tilde{\nu})$ 的最优计划 $\hat{\pi}$。
    \item 仍令 $\hat{\pi}_{\mathrm{new}}=(\pi-\pi')+Z\hat{\pi}$。上一帧已验证：它是 $(\mu,\nu)$ 的竞争计划。
    \item 因 $\hat{\pi}$ 与 $\tilde{\pi}$ 都最优，它们成本相同，故
  \end{itemize}
  \[
    I[\hat{\pi}_{\mathrm{new}}]
    = I[\pi]-Z I[\tilde{\pi}]+Z I[\hat{\pi}]
    = I[\pi].
  \]
  所以 $\hat{\pi}_{\mathrm{new}}$ 也是 $(\mu,\nu)$ 的最优计划；由 $\pi$ 的唯一性，
  $\hat{\pi}_{\mathrm{new}}=\pi$。消去公共部分得 $Z\hat{\pi}=\pi'=Z\tilde{\pi}$，
  又 $Z>0$，故 $\hat{\pi}=\tilde{\pi}$。\hfill$\square$
  \textbf{例 4.7（条件化仍最优）。}
  设 $(X,Y)$ 是最优计划 $\pi$ 下的随机向量，$Z\subset\X\times\Y$ Borel 且 $\Prob[(X,Y)\in Z]>0$。
  令 $\pi' = \pi|_{Z}$，则 $\pi'\le\pi$，
  定理 4.6 直接给出：给定 $(X,Y)\in Z$ 的条件耦合在相应边际之间最优。
  \TLtakeaway{唯一性传递的关键等式是 $I[\hat{\pi}_{\mathrm{new}}]=I[\pi]$：整体唯一性迫使局部替代等于原局部计划。}
\end{frame}

% ── 例 4.7 的原则解释 ──────────────────────────────────────────────────────────
\begin{frame}{例 4.7：条件化原则的威力}
  % Manual warn-style block (avoids 0.4pt draw-border overflow from \TLwarnblock)
  \begin{tikzpicture}
    \node[draw=TLneg, fill=TLneg!8, rounded corners=3pt,
          inner sep=8pt, text width=\linewidth-17pt, align=justify] {%
      \small\textbf{\color{TLneg}为什么这个原则"简单但有力"}\\[0.3em]
      全书多次用到"取一个最优计划，把注意力限制到某个子集，
      子集上的计划仍最优"这一论证步骤。每次背后都是定理 4.6。
    };
  \end{tikzpicture}

  \textbf{使用场景举例：}
  \begin{itemize}
    \item 把最优计划限制到某个"好的"区域，在该区域上研究其几何结构。
    \item 用条件化把全局问题降维为局部问题，再用局部最优性归纳推进。
    \item 在唯一性证明中：若存在两个最优计划，对它们的"差集"条件化推出矛盾。
  \end{itemize}

  \TLtakeaway{例 4.7：最优耦合对任意正概率事件的条件化也最优（定理 4.6 之概率语言形式）。}
\end{frame}

% ── 定理 4.8：最优成本的凸性 ──────────────────────────────────────────────────
\begin{frame}{第三件好事：最优成本的凸性（定理 4.8）}
  \deflab{定理 4.8（最优成本的凸性，Villani Thm 4.8）。}
  设 $(\Theta,\lambda)$ 为概率空间，$\theta\mapsto\mu_{\theta}\in\Pp(\X)$，
  $\theta\mapsto\nu_{\theta}\in\Pp(\Y)$ 可测，成本 $c\ge a+b$ 如定理 4.1。记混合边际
  $\mu=\int_{\Theta}\mu_{\theta}\dd\lambda(\theta)$，
  $\nu=\int_{\Theta}\nu_{\theta}\dd\lambda(\theta)$。则
  \[
    C(\mu,\nu)
    \;\le\;
    \int_{\Theta} C(\mu_{\theta},\nu_{\theta})\dd\lambda(\theta).
  \]

  \textbf{直觉。} "先混合，再最优传输"的代价不超过"先分别最优传输，再混合"的平均代价。
  混合只会\emph{增加灵活性}（可以选混合后的最优计划），故成本不增。

  \deflab{读法。} 这不是说每个最优计划线性变化；它只比较\emph{最优值}。右端给出一个可行的混合方案，
  左端允许在混合后的更大可行集中重新优化。

  \TLtakeaway{最优传输成本关于边际测度\emph{凸}：混合分布的传输成本不超过分别传输成本的平均。}
\end{frame}

\begin{frame}{定理 4.8 的证明（一）：混合计划有正确边际}
  \textbf{选取局部最优计划。}\proofskipnote 对 $\lambda$-a.e.\ $\theta$，定理 4.1 给出
  $\pi_{\theta}\in\coup{\mu_{\theta}}{\nu_{\theta}}$ 且
  $I[\pi_{\theta}]=C(\mu_{\theta},\nu_{\theta})$。
  Villani 暂时借用推论 5.22，保证可把 $\theta\mapsto\pi_{\theta}$ 选成可测。

  \deflab{补一个事实：测度值积分的边际。}
  令 $\pi=\int_{\Theta}\pi_{\theta}\dd\lambda(\theta)$。任取 $\varphi\in C_b(\X)$，把边际推出积分号：
  \[
    \begin{aligned}
      \int_{\X}\varphi\dd\pf{(\mathrm{proj}_{\X})}{\pi}
      &= \int_{\Theta}\!\int_{\X}\varphi\dd\pf{(\mathrm{proj}_{\X})}{\pi_{\theta}}\dd\lambda(\theta)\\
      &= \int_{\Theta}\!\int_{\X}\varphi\dd\mu_{\theta}\dd\lambda(\theta)
       = \int_{\X}\varphi\dd\mu .
    \end{aligned}
  \]
  故第一边际为 $\mu$；同理第二边际为 $\nu$，所以 $\pi\in\coup{\mu}{\nu}$。

  \TLtakeaway{边际核对要用检验函数：把 $\mathrm{proj}_{\X}$ 的边际积分推过 $\theta$-积分，得到混合后的第一边际 $\mu$；第二边际同理。}
\end{frame}

\begin{frame}{定理 4.8 的证明（二）：Fubini 给出成本上界}
  \textbf{用上帧的可行计划估计最优值。}\proofskipnote 已知
  $\pi=\int_{\Theta}\pi_{\theta}\dd\lambda(\theta)\in\coup{\mu}{\nu}$，所以
  $C(\mu,\nu)\le I[\pi]$。

  \deflab{补一个事实：Fubini 定理。}
  在本处它允许交换参数 $\theta$ 与位置 $(x,y)$ 的积分次序，把混合计划的总成本拆成局部计划成本的平均。

  \[
    \begin{aligned}
      C(\mu,\nu)
      &\le I[\pi]
       = \int_{\X\times\Y} c\dd\pi\\
      &= \int_{\Theta}\!\left(\int_{\X\times\Y} c\dd\pi_{\theta}\right)\dd\lambda(\theta)
       = \int_{\Theta} C(\mu_{\theta},\nu_{\theta})\dd\lambda(\theta).
    \end{aligned}
  \]
  最后一行用 $I[\pi_{\theta}]=C(\mu_{\theta},\nu_{\theta})$，这是每个 $\pi_{\theta}$ 的最优性。
  \TLtakeaway{成本上界只需一个可行计划：混合局部最优计划，再用 Fubini 把其成本拆回局部最优成本的平均。}
\end{frame}

% ── 最优计划的刻画：开放问题（一）───────────────────────────────────────────
\begin{frame}{超越存在性：最优计划的刻画（一）}
  存在性定理告诉我们"最优计划存在"，但对其\emph{结构}只字未提。自然的问题：

  \begin{itemize}
    \item \textbf{唯一性？} 最优计划是否唯一？是否光滑？
    \item \textbf{Monge 耦合？} 是否存在\emph{确定性}最优计划（即推前映射 $T$ 使 $\pf{T}{\mu}=\nu$）？
    \item \textbf{几何刻画？} 有没有直接判断某个耦合"是否最优"的几何条件？
  \end{itemize}

  关于唯一性与几何刻画的答案见第 9--12 章；关于几何刻画的关键工具——循环单调性与 Kantorovich 对偶——将在下一节（§2）展开。
\end{frame}

% ── 最优计划的刻画：开放问题（二）───────────────────────────────────────────
\begin{frame}{超越存在性：为何 Monge 不能用紧性论证（二）}
  \textbf{为什么不能直接对 Monge 用定理 4.1 的紧性论证？}

  Monge 计划的集合（由推前映射产生的耦合）\emph{一般不弱紧}——事实上，它往往在全体耦合 $\coup{\mu}{\nu}$ 中\emph{稠密}！
  因此：
  \begin{itemize}
    \item Monge 问题的下确界 $=$ Kantorovich 问题的最小值（两者相等是\emph{可能}的）。
    \item 但\emph{无法保证} Monge 极小化子存在——下确界可以不被 Monge 计划取到。
  \end{itemize}

  \textbf{对比：}
  \begin{itemize}\small
    \item Kantorovich 松弛允许"分裂质量"，集合 $\coup{\mu}{\nu}$ 弱紧，极小化子必然存在。
    \item Monge 限制不允许分裂，集合不紧，极小化子可以不存在。
  \end{itemize}

  \TLtakeaway{确定性 (Monge) 计划集不紧，因此 Monge 极小子不保证存在；但 Kantorovich 最小值 $=$ Monge 下确界仍可能成立。第 5 章开始刻画。}
\end{frame}

% ── 例 4.9：质量分裂（TikZ 图 4.1）── 文字 ──────────────────────────────────
\begin{frame}{例 4.9：需要分裂质量的例子}
  \deflab{例 4.9（Villani Example 4.9）。}
  取 $\X=\Y=\R^2$，成本 $c(x,y)=\abs{x-y}^2$。
  设：
  \begin{align*}
    \mu &= H^1 \text{ 限制到 } \{0\}\times[-1,1], \\
    \nu &= \tfrac{1}{2}H^1 \text{ 限制到 } \{-1\}\times[-1,1]
           + \tfrac{1}{2}H^1 \text{ 限制到 } \{1\}\times[-1,1].
  \end{align*}
  其中 $H^1$ 为一维 Hausdorff 测度（弧长测度）。

  \textbf{唯一最优计划：} 对每个点 $(0,a)\in\supp\mu$，将其一半质量送到 $(-1,a)$，
  另一半送到 $(1,a)$（水平运输，竖直坐标不变）：
  \[
    \pi^\star = \int_{-1}^{1} \tfrac{1}{2}\bigl[\delta_{((0,a),(-1,a))}+\delta_{((0,a),(1,a))}\bigr]\dd a.
  \]

  \textbf{这\emph{不是} Monge 计划：} 每个 $(0,a)$ 的质量被分裂成两份，
  不存在函数 $T$ 使 $T(0,a)=(-1,a)$ 且同时 $T(0,a)=(1,a)$。

  \TLtakeaway{质量分裂是最优的——这说明即使有唯一最优 Kantorovich 计划，也不一定是 Monge 计划。}
\end{frame}

% ── 例 4.9：TikZ 图 4.1 ────────────────────────────────────────────────────
\begin{frame}{例 4.9 图示：最优耦合 vs 确定性近似（图 4.1 重绘）}
  \begin{columns}[T]
    \begin{column}{0.5\textwidth}
      \centering{\small\textbf{最优耦合（质量分裂）}}\\[6pt]
      \begin{tikzpicture}[scale=1.05,>={Stealth[length=4pt]}]
        % 支撑集
        \draw[TLprimary,very thick] (0,-1.3) -- (0,1.3);
        \draw[TLaccent,thick] (-1.5,-1.3) -- (-1.5,1.3);
        \draw[TLaccent,thick] (1.5,-1.3) -- (1.5,1.3);
        % 箭头：中心 -> 左和右（几个示例点）
        \foreach \y in {-1.0,-0.5,0.0,0.5,1.0}{
          \draw[->,TLink,thin] (0.05,\y) -- (1.45,\y);
          \draw[->,TLink,thin] (-0.05,\y) -- (-1.45,\y);
        }
        % 标注
        \node[font=\scriptsize,TLinkSoft,above] at (0,1.35) {$\mu$};
        \node[font=\scriptsize,TLinkSoft,above] at (-1.5,1.35) {$\nu_L$};
        \node[font=\scriptsize,TLinkSoft,above] at (1.5,1.35) {$\nu_R$};
        \node[font=\scriptsize] at (0,-1.6) {$x=0$};
        \node[font=\scriptsize] at (-1.5,-1.6) {$x=-1$};
        \node[font=\scriptsize] at (1.5,-1.6) {$x=1$};
      \end{tikzpicture}\\[3pt]
      {\scriptsize 每点质量各分一半，水平运输。}
    \end{column}
    \begin{column}{0.5\textwidth}
      \centering{\small\textbf{确定性近似（非最优）}}\\[6pt]
      \begin{tikzpicture}[scale=1.05,>={Stealth[length=4pt]}]
        \draw[TLprimary,very thick] (0,-1.3) -- (0,1.3);
        \draw[TLaccent,thick] (-1.5,-1.3) -- (-1.5,1.3);
        \draw[TLaccent,thick] (1.5,-1.3) -- (1.5,1.3);
        % 交替箭头：显式指定每个点的去向
        \draw[->,TLink,thin] (0.05,-1.0) -- (1.45,-1.0);
        \draw[->,TLink,thin] (-0.05,-0.5) -- (-1.45,-0.5);
        \draw[->,TLink,thin] (0.05,0.0) -- (1.45,0.0);
        \draw[->,TLink,thin] (-0.05,0.5) -- (-1.45,0.5);
        \draw[->,TLink,thin] (0.05,1.0) -- (1.45,1.0);
        \node[font=\scriptsize,TLinkSoft,above] at (0,1.35) {$\mu$};
        \node[font=\scriptsize,TLinkSoft,above] at (-1.5,1.35) {$\nu_L$};
        \node[font=\scriptsize,TLinkSoft,above] at (1.5,1.35) {$\nu_R$};
        \node[font=\scriptsize] at (0,-1.6) {$x=0$};
        \node[font=\scriptsize] at (-1.5,-1.6) {$x=-1$};
        \node[font=\scriptsize] at (1.5,-1.6) {$x=1$};
      \end{tikzpicture}\\[3pt]
      {\scriptsize 每点只送一侧（交替），不分裂，非最优。}
    \end{column}
  \end{columns}
  \vspace{4pt}
  {\scriptsize 据 Villani 图 4.1 重绘。}
  \TLtakeaway{最优计划（左）把质量对称分裂；确定性近似（右）不分裂但成本更高。Monge 下确界 $=$ Kantorovich 最小值，但 Monge 极小子不存在。}
\end{frame}

% ── 习题 1 · Exercises ─────────────────────────────────────────────────────────
\begin{frame}{习题 1 · Exercises}
  \begin{itemize}
    \item \textbf{习题 1.1}~(\diffI)\ 验证独立耦合 $\mu\otimes\nu$ 满足 Kantorovich 问题的边际约束，
      并在 $\iint c\dd\mu\dd\nu<+\infty$ 时说明 $I[\mu\otimes\nu]<+\infty$，
      从而 $C(\mu,\nu)<+\infty$。
      \begin{itemize}\footnotesize
        \item \hint{验证 $\pf{(\mathrm{proj}_{\X})}{\mu\otimes\nu}=\mu$ 和 $\pf{(\mathrm{proj}_{\Y})}{\mu\otimes\nu}=\nu$，
          再利用 Fubini 定理计算 $I[\mu\otimes\nu]=\iint c\dd\mu\dd\nu$。}
      \end{itemize}
    % SOLUTION: mu otimes nu 的第一边际：对任意 Borel A 属于 X，mu otimes nu [A × Y] = mu[A] * nu[Y] = mu[A]，故第一边际为 mu；同理第二边际为 nu。I[mu otimes nu] = int_{X×Y} c d(mu⊗nu) = int_X int_Y c(x,y) dnu(y) dmu(x) (Fubini，c>=0 可以交换) = ∬ c dmu dnu < +∞ 由假设。故独立耦合有限成本，C(μ,ν) ≤ I[μ⊗ν] < +∞。

    \item \textbf{习题 1.2}~(\diffII)\ 设 $\pi^\star$ 是有限离散分布的最优 Kantorovich 计划。
      取子集 $Z=\{(a_i,b_j):\pi^\star_{ij}>0\}$ 中某个矩形子块 $A\times B$，
      令 $\pi'=\pi^\star|_{A\times B}$。用显式"重新分配"论证证明 $\pi'/\pi'[A\times B]$ 是最优的。
      \begin{itemize}\footnotesize
        \item \hint{假设子计划不最优，构造更便宜的局部替换，再代入整体计划，与 $\pi^\star$ 最优性矛盾。}
      \end{itemize}
    % SOLUTION: 设子计划 tilde{pi} = pi'/(pi'[A×B]) 不是最优的。则存在有相同边际但成本更小的计划 hat{pi}。令 Z = pi'[A×B]，构造 pi_new = (pi^star - pi') + Z * hat{pi}。验证：pi_new 的边际等于 mu, nu（因为 pi' 和 hat{pi} 有相同的边际，差了一个因子 Z）；pi_new 的成本 = I[pi^star] - I[pi'] + Z * I[hat{pi}] = I[pi^star] - Z * I[tilde{pi}] + Z * I[hat{pi}] < I[pi^star]（因为 I[hat{pi}] < I[tilde{pi}]）。这与 pi^star 最优性矛盾。

    \item \textbf{习题 1.3}~(\diffIII)\ 用例 4.9 说明：存在 $\mu,\nu\in\Pp(\R^2)$ 使得
      Monge 下确界等于 Kantorovich 最小值，但不存在 Monge 映射 $T$ 实现 Kantorovich 最小值。
      \begin{itemize}\footnotesize
        \item \hint{例 4.9 中唯一最优 Kantorovich 计划分裂质量，不对应任何映射；
          但确定性近似序列的成本趋向 $C(\mu,\nu)$，说明两个下确界相等。}
      \end{itemize}
    % SOLUTION: 例 4.9 中 C(mu,nu) 由唯一最优计划 pi^star 达到，其值为 1（每点 (0,a) 到 (±1,a) 的平方距离均为 1，总成本为 int_{-1}^{1} 1 da = 2，但归一化后为 1）。任意 Monge 计划 T 满足 pf{T}{mu} = nu，必将 [−1,1] 的一半质量送到 x=−1 段，另一半送到 x=1 段；设在集合 A 上 T(0,a) = (−1,a)，在 [−1,1]\A 上 T(0,a)=(1,a)，A 测度为 1。任意这样的 T 成本等于 1 = C(mu,nu)，但没有一个是连续映射或弱极限——因为 pi^star 不是推前测度。实际上可以验证 pi^star 的条件分布 pi^star_{(0,a)} = 1/2 delta_{(-1,a)} + 1/2 delta_{(1,a)} 不是 delta 测度，故 pi^star 不是任何映射的图。但 Monge 下确界确实等于 Kantorovich 最小值（均为 1），因为上述显式 Monge 映射（取 A 任意测度为 1 的子集）成本恰好也是 1。此例说明：Kantorovich 最小值可以被 Monge 计划的成本序列实现，但极限（唯一最优 pi^star）本身不是 Monge 计划。
  \end{itemize}
\end{frame}

% ── §1 小结 ───────────────────────────────────────────────────────────────────
\begin{frame}{§1 小结：最优耦合的四个基本性质}
  \begin{itemize}\small
    \item \textbf{存在性（定理 4.1）。} Polish 空间上下半连续成本满足下界条件时，最优耦合必然存在。
      证明骨架：Prokhorov 紧性 $+$ 成本泛函下半连续。
    \item \textbf{成本有定义（注记 4.2 / 4.5）。} 下界条件保证 $I[\pi]$ 在 $\R\cup\{+\infty\}$ 有意义；
      独立耦合成本有限 $\iff$ $\iint c\dd\mu\dd\nu<+\infty$。
    \item \textbf{限制性质（定理 4.6）。} 最优计划的任意子计划也最优——"局部也优"。
      条件化同样保持最优性。
    \item \textbf{凸性（定理 4.8）。} 最优成本 $C(\mu,\nu)$ 关于边际测度凸：
      混合分布的传输成本不超过分别传输成本的平均。
  \end{itemize}
  \textbf{遗留问题（§2 起逐步回答）：}
  \begin{itemize}\small
    \item 唯一性？几何刻画？$\to$ §2 循环单调性与 Kantorovich 对偶。
    \item 确定性（Monge）最优计划是否存在？$\to$ 第 9--12 章。
  \end{itemize}
  \TLtakeaway{第 4 章打下了基础：存在性、有限性、限制性、凸性——四块基石，为后续刻画最优计划做好准备。}
\end{frame}
